\centerline{\bf \Huge Теория Чисел. База}

Делимость

Кратное число - это число, которое делится на другое число без остатка. 

Делитель - это число, на которое делится другое число без остатка

Делимое, делитель, частное, неполное частное, остаток (0...m-1). (1-2 задачи)

НОД, НОК. Алгоритм вычисления. (3 задача)

Полезно знать: Факториал. Полный квадрат.

(4,5 задачи)

Поиск всех делителей числа. Руками и вывод формулы.

\problem{\textbf{1}}  Два класса купили 737 учебников. Каждому досталось одно и то же число учебников. Сколько было ребят в обоих классах вместе?

\problem{\textbf{2}} Делимое равно 371, остаток 30. Найти делитель и соответствующее ему неполное частное.

\problem{\textbf{3}} Найдите НОД и НОК следующих чисел:

а) 36 и 54;

б) 50 и 81;

в) 2520 и 13200.

\problem{\textbf{4}} Перечислите все делители чисел: 

а) 3·5·5·11;

б) 5·45; 

в) 1001;

г) 256.

\problem{\textbf{5}} 
Сколько натуральных делителей имеет число, если \(p\) и \(q\) – простые числа: 

а) \(pq^2\);

б) \(p^2q^3\);

в) \(p^4q^5\);

г) \(p^q\).


